\section{Satellites}
A satellite is an object that \red{orbits another object}.\\

There are \textbf{natural} satellites and \textbf{artificial} satellites.
\begin{itemize}
    \item The Moon is a \textbf{natural} satellite of the Earth.
    \item The International Space Station (ISS) is an \textbf{artificial} satellite.
\end{itemize}

\subsection{Newton's Cannon}
His idea was: \textit{if a cannon is placed on top of a very tall mountain and air resistance is ignored, the cannon 
shoots a cannonball horizontally}.\\

At the ideal speed, the distance the cannonball falls \textbf{equals} the distance that the Earth's surface curves away beneath it.\\

If $v < v_{\text{ideal}}$, the distance between the cannonball and Earth's surface will \textbf{decrease}.\\

If $v > v_{\text{ideal}}$, the distance between the cannonball and Earth's surface will \textbf{increase}.\\

The cannonball must have a constant speed and travel along a perfect circular path.

\subsection{Geosynchronous}
They have the same orbital period as the \textbf{rotation period} of the object they orbit.\\

The period is around \textbf{24 hrs}\\

There is a special type of geosynchronous orbit called \textbf{geostationary}.
\begin{itemize}
    \item Remain above a location on Earth's \textbf{equator}
    \item Orbit in the same \textbf{direction} that the Earth rotates
\end{itemize}

\subsection{Formulas related to satellite}
We will derive each formula in this handout.\\

To begin, let us draw the FBD for the satellite.\\
\begin{center}
    \begin{tikzpicture}[scale=1.2, >=Stealth]

    % Earth
    \shade[ball color=blue!40] (0,0) circle (1);
    \node at (0,0) {\textbf{Earth}};

    % Satellite position
    \coordinate (satellite) at (0,3);
    \filldraw[gray!60] (satellite) circle (0.15);
    \node[above=2pt of satellite] {Satellite};

    % Forces
    \draw[thick, red, ->] (satellite) -- ++(0,-1.2)
        node[midway, right] {$F_g$};
    \draw[thick, blue, ->] (satellite) -- ++(0,1.2)
        node[midway, right] {$F_{\text{fict}}$};

    % Orbit line
    \draw[dashed, gray] (0,0) circle (3);
    \node[right] at (0,1.5) {Orbit altitude};

    % Optional: radius vector
    \draw[->, thick, black!50] (0,0) -- (satellite)
        node[midway, left] {$r$};

    \end{tikzpicture}
\end{center}

\textbf{Down} is negative\\

Let us derive the formula for a satellite:
\begin{align*}
    \sum \vec{F} &= m\vec{a_{per}}\\
    F_g - F_{fict} &= 0\\
    mg - m\left|F_{FOR}\right| &= 0\\
    mg &= ma_c\\
    a_c &= \frac{Gm}{R^2} 
\end{align*}

Sub in $a_c = \frac{V^2}{R}$:
\begin{align*}
    \frac{v^2}{R} &= \frac{GM}{R^2}\\
    v &= \sqrt{\frac{GM}{R}}
\end{align*}

\begin{center}
    $v$ = orbital speed\\
    $M$ = mass of the object being orbited (kg)
\end{center}

Sub in $a_c = \frac{4\pi^2 R}{T^2}$
\begin{align*}
    \frac{4\pi^2 R}{T^2} &= \frac{GM}{R}\\
    T^2 &= \frac{4\pi^2 R^3}{GM}\\
    T &= \sqrt{\frac{4\pi^2 R^3}{GM}}
\end{align*}
